
\documentclass[11pt]{article}

\usepackage[margin=1in]{geometry}
\usepackage[T1]{fontenc}
\usepackage{lmodern}
\usepackage{microtype}
\usepackage{amsmath,amssymb,mathtools}
\usepackage{booktabs,longtable,array}
\usepackage{enumitem}
\usepackage{xcolor}
\usepackage{listings}
\usepackage[hidelinks]{hyperref}
\usepackage{fancyhdr}
\usepackage{titlesec}

\definecolor{codegray}{gray}{0.96}
\definecolor{darkgray}{gray}{0.25}

\lstset{
  basicstyle=\ttfamily\small,
  backgroundcolor=\color{codegray},
  frame=single,
  rulecolor=\color{gray!35},
  breaklines=true,
  columns=fullflexible,
  keepspaces=true,
  showstringspaces=false,
  xleftmargin=0.5em,
  xrightmargin=0.5em
}

\pagestyle{fancy}
\fancyhf{}
\lhead{\texttt{netcenter} v0.1.0}
\rhead{Algorithms and implementation}
\cfoot{\thepage}
\setlength{\headheight}{14pt}

\titleformat{\section}{\Large\bfseries}{\thesection}{0.75em}{}
\titleformat{\subsection}{\large\bfseries}{\thesubsection}{0.75em}{}

\newcommand{\R}{\mathbb{R}}
\newcommand{\argmin}{\operatorname*{arg\,min}}
\newcommand{\argmax}{\operatorname*{arg\,max}}
\newcommand{\ecc}{\operatorname{ecc}}
\newcommand{\clip}{\operatorname{clip}}

\title{
\textbf{Technical Discussion of \texttt{netcenter}}\\
\large Network Construction, Shortest Paths, the 1-Median, and the Absolute 1-Center
}

\author{Technical note for \texttt{netcenter} v0.1.0}
\date{24 September 2026}

\begin{document}

\maketitle

\begin{abstract}

\texttt{netcenter} is a Python package for solving classical facility-location
problems on undirected spatial networks.

Given a road or other line-based network and an optional set of demand
locations, the package computes:

\begin{enumerate}[leftmargin=1.5em]
\item a weighted network 1-median;
\item a vertex-restricted 1-center; and
\item an exact absolute 1-center whose optimum may lie in the interior of a
network edge.
\end{enumerate}

Distances are shortest-path distances on the constructed network rather than
Euclidean distances in the plane. The facility-location mathematics is
classical: the median and absolute-center formulation follows Hakimi, while the
continuous center computation follows the edgewise structure developed in the
network-location literature, particularly Kariv and Hakimi. The package does
not claim a new location algorithm.

The package-specific work lies primarily in geospatial network construction,
topology safeguards, numerical handling, memory management, vectorized
implementation, pruning, parallel execution, and regression testing. This
document describes both the mathematics and the implementation provenance of
the first public release, v0.1.0, including the main pre-release corrections and
engineering decisions.

\end{abstract}

\section{Scope and notation}

Let

\[
G=(V,E)
\]

be a connected undirected network with non-negative edge lengths.

Let

\[
Q=\{q_1,\ldots,q_k\}\subseteq V
\]

be the unique demand vertices after external observations have been mapped to
the network and duplicate locations have been coalesced.

Let

\[
d_G(x,y)
\]

denote shortest-path distance in the network.

Write

\begin{itemize}[leftmargin=1.5em]
\item \(n=|V|\) for the number of network vertices;
\item \(m=|E|\) for the number of network edges;
\item \(k=|Q|\) for the number of unique demand vertices;
\item \(L\) for the length of a particular edge.
\end{itemize}

The package currently solves three related location problems.

\subsection{Weighted 1-median}

Given non-negative demand weights \(w_i\),

\[
x^*
=
\argmin_{x\in V}
\sum_{i=1}^{k}
w_i d_G(x,q_i).
\]

This is the network minisum objective.

For vertex demand, Hakimi's vertex-optimality result implies that an optimal
median may be chosen at a network vertex \cite{Hakimi1964}. The implementation
therefore evaluates network vertices rather than searching edge interiors.

\subsection{Vertex 1-center}

The vertex 1-center solves

\[
x^*
=
\argmin_{x\in V}
\max_i d_G(x,q_i).
\]

This is a node-restricted minimax objective.

\subsection{Absolute 1-center}

The absolute 1-center allows the candidate location to lie anywhere on the
network:

\[
x^*
=
\argmin_{x\in G}
\max_i d_G(x,q_i).
\]

Unlike the median, an optimal absolute center need not occur at a vertex and may
lie strictly inside an edge \cite{Hakimi1964}. This edge-interior case is the
principal mathematical reason for the package.

\section{Algorithmic lineage}

Every major mathematical component of \texttt{netcenter} is classical.

The principal references are:

\begin{itemize}[leftmargin=1.5em]

\item Hakimi \cite{Hakimi1964}, who formulated the absolute center and median
of a graph and established the key vertex-optimality result used by the median
solver;

\item Hakimi \cite{Hakimi1965}, who extended the network-location framework to
multiple facilities;

\item Kariv and Hakimi \cite{KarivHakimi1979a}, who developed algorithmic
methods for network \(p\)-center problems, including the continuous
absolute-center case;

\item Kariv and Hakimi \cite{KarivHakimi1979b}, for the corresponding
\(p\)-median literature;

\item Handler and Mirchandani \cite{HandlerMirchandani1979}, for a systematic
treatment of location on networks and the classical edgewise center
formulation;

\item Daskin \cite{Daskin2013}, for a modern treatment of network and discrete
facility-location models.

\end{itemize}

Additional algorithms used as supporting infrastructure are also classical:

\begin{itemize}[leftmargin=1.5em]

\item Dijkstra's shortest-path algorithm \cite{Dijkstra1959};

\item k-d trees for nearest-neighbor lookup, due to Bentley
\cite{Bentley1975};

\item connected-component analysis using standard graph traversal, with the
classical algorithmic lineage represented by Hopcroft and Tarjan
\cite{HopcroftTarjan1973};

\item sweep-line segment-intersection algorithms such as Bentley--Ottmann
\cite{BentleyOttmann1979}, relevant to planar noding;

\item finite-precision geometric techniques including Greene and Yao
\cite{GreeneYao1986} and Hobby \cite{Hobby1999}, relevant to robust geometric
intersection and snap rounding.

\end{itemize}

The package does not reimplement all of these algorithms. Several are delegated
to SciPy, Shapely, GEOS, and related numerical/geospatial libraries.

\section{From spatial linework to a graph}

A GIS line layer records geometry, not necessarily routing topology.

Two line geometries may cross in the plane without permitting movement between
them. Conversely, a genuine junction may occur at a coordinate appearing in
the interior of one or more LineStrings.

Accordingly, graph construction is part of the model, not merely a
file-conversion step.

\subsection{Default topology rule: shared-source-vertex noding}

The default behavior is \emph{shared-source-vertex noding}.

If a snapped coordinate already appears as a source vertex in at least two
distinct input LineStrings, the location is treated as a junction and the
participating lines are split there.

This recovers common T-junctions and shared interior-interior junctions already
encoded in source linework.

It does not create a new junction merely because two segments cross
geometrically in two dimensions.

This distinction reduces the risk of false turns at:

\begin{itemize}[leftmargin=1.5em]
\item bridges;
\item flyovers;
\item tunnels;
\item grade-separated intersections.
\end{itemize}

The shared-source-vertex rule is straightforward topology recovery rather than
a new graph algorithm.

\subsection{Endpoint-only mode}

If the source is already segmented at every true junction, the caller may
disable shared-vertex splitting.

In this mode, connectivity is inferred strictly from line endpoints.

\subsection{Explicit planar noding}

The caller may explicitly request full planar noding.

In that mode, geometric crossings are converted into junctions.

This is appropriate only when every two-dimensional crossing genuinely
represents connectivity.

The sweep-line method for reporting all intersections among a set of segments
is due to Bentley and Ottmann \cite{BentleyOttmann1979}. Practical finite
precision intersection handling and snap rounding have their own computational
geometry literature, including Greene and Yao \cite{GreeneYao1986} and Hobby
\cite{Hobby1999}.

\texttt{netcenter} does not implement those algorithms directly. Planar noding
is delegated to Shapely/GEOS \cite{ShapelyNode}. GEOS itself descends from the
JTS Topology Suite.

\subsection{Shape vertices and topology vertices}

Road geometries often contain many intermediate coordinates whose purpose is
only to describe curvature.

Those coordinates are geometrically useful but need not all survive as graph
vertices.

Reducing unnecessary shape vertices matters because the shortest-path matrix
scales with the number of graph nodes.

Line merging is delegated to Shapely/GEOS.

\subsection{Matching near-identical coordinates}

Endpoints intended to coincide may differ at very small floating-point scales.

The graph builder quantizes coordinates to a configurable grid before assigning
node identity. This is an engineering convention rather than a named algorithm.

\subsection{Parallel edges}

Real road data may contain multiple segments joining the same pair of graph
vertices.

For shortest-path purposes, the sparse adjacency stores the minimum direct edge
length between a node pair.

This is important because converting a sparse coordinate matrix to CSR can sum
duplicate entries if not handled carefully. Summing parallel road lengths would
be wrong and can violate the metric assumptions used by the center sweep.

\subsection{Closed and nearly closed rings}

A closed LineString has coincident start and end coordinates.

If treated naively, it can become a self-loop and disappear during graph
cleaning.

The graph builder detects closed and nearly closed rings and splits them before
graph construction so that the physical network remains represented.

\subsection{Filtered slivers and orphan nodes}

Small slivers or snapped self-loops may be removed after preliminary node IDs
have been assigned.

The current implementation compacts node IDs after edge filtering so removed
edges do not leave isolated matrix rows or columns that later produce spurious
infinite distances.

\subsection{Connected components}

Real line datasets often contain isolated fragments.

Identifying connected components is standard graph traversal; the classical
algorithmic lineage includes Hopcroft and Tarjan \cite{HopcroftTarjan1973}.
The package uses SciPy's sparse graph routines rather than implementing a
component algorithm directly.

When component filtering is enabled, the retained component is selected by
total road length rather than merely node count.

\section{Coordinate systems and metric units}

The mathematical optimization assumes meaningful edge lengths.

The package therefore operates in a projected coordinate reference system whose
units are metres.

The policy is:

\begin{itemize}[leftmargin=1.5em]

\item projected metre-based input is preserved;

\item geographic input is transformed to an inferred local UTM CRS unless an
explicit target CRS is supplied;

\item projected input using non-metre units is transformed;

\item an explicitly supplied target CRS must be projected and metre-based.

\end{itemize}

For geographically large study regions, the analyst should supply a projection
appropriate to the full study area rather than rely automatically on a single
UTM zone.

\section{Demand representation}

External demand observations need not coincide exactly with graph vertices.

The current package maps each demand observation to the nearest network node.

For a point observation \(p_i\),

\[
q_i
=
\argmin_{v\in V}
\|p_i-v\|_2.
\]

The nearest-node lookup is performed with a k-d tree, using SciPy's
\texttt{cKDTree}; the k-d tree was introduced by Bentley
\cite{Bentley1975}.

A maximum snapping distance may be imposed.

This is an explicit modeling assumption: demand is represented at vertices, not
continuously along edges.

\subsection{Non-point geometries}

If a demand layer contains polygons or other non-point geometries, each usable
geometry is converted to a representative point before snapping.

\subsection{Repeated snapped demand}

Several observations may snap to the same network node.

For the minimax objectives, multiplicity has no effect:

\[
\max\{d(x,v),d(x,v)\}
=
d(x,v).
\]

Duplicate center demands can therefore be removed exactly.

For the median, multiplicity matters. Their weights are summed:

\[
w_v
=
\sum_{i:q_i=v} w_i.
\]

Thus duplicate-demand coalescing reduces shortest-path work without changing
either optimization problem.

\section{Shortest-path matrix}

After demand coalescing, the package computes a shortest-path matrix

\[
D\in\R^{k\times n},
\]

where

\[
D_{iv}
=
d_G(q_i,v).
\]

Shortest paths are computed using Dijkstra's algorithm
\cite{Dijkstra1959} through SciPy's sparse graph implementation
\cite{SciPyDijkstra}.

The graph is required to be square, non-negative, and symmetric because the
package currently models undirected networks.

\subsection{Computational cost}

For \(k\) demand sources, \(n\) nodes, and \(m\) edges, a conventional
binary-heap complexity statement is approximately

\[
O\!\left(k(m+n\log n)\right),
\]

although actual SciPy performance depends on implementation details and graph
structure.

The matrix itself requires

\[
O(kn)
\]

storage.

In many practical instances, memory usage of \(D\) is a more important limit
than the arithmetic required by the location objective.

\subsection{Chunking}

Dijkstra calculations are performed in bounded groups of source nodes.

If a block contains \(r\) demand sources, its temporary output contains
approximately \(rn\) floating-point values.

The \texttt{max\_temp\_mb} parameter controls the target size of one temporary
float64 block per worker.

\subsection{Streaming parallel assembly}

When several Dijkstra blocks are computed in parallel, completed blocks are
written directly into a preallocated final matrix.

This avoids retaining every completed block and then constructing another full
matrix by concatenation.

\subsection{Reduced precision}

The final shortest-path matrix may optionally be stored in float32.

Approximate storage requirements are:

\[
8kn \text{ bytes}
\]

for float64 and

\[
4kn \text{ bytes}
\]

for float32.

Float32 is a memory--precision trade-off. The optimization remains internally
consistent with the stored matrix, but the stored metric is itself rounded
relative to float64 output.

\section{Weighted 1-median}

For every candidate node \(v\), define

\[
M(v)
=
\sum_{i=1}^{k}
w_iD_{iv}.
\]

The weighted network median is

\[
v^*
=
\argmin_{v\in V}M(v).
\]

Hakimi's vertex-optimality result reduces the continuous network median problem
to this finite evaluation for vertex demand \cite{Hakimi1964}.

In matrix notation, if

\[
w=
\begin{bmatrix}
w_1 & \cdots & w_k
\end{bmatrix},
\]

then

\[
M=wD.
\]

The implementation is therefore a weighted reduction over the shortest-path
matrix followed by an \texttt{argmin}.

If no explicit weights are supplied, unit weights are used.

With unit weights, minimizing the total distance is closely related to
maximizing closeness centrality, whose classical lineage includes Bavelas
\cite{Bavelas1950} and Freeman \cite{Freeman1978}.

On trees, related median and center problems admit simpler algorithms; Goldman
\cite{Goldman1971} is a classical reference.

\section{Vertex 1-center}

For each candidate vertex \(v\), define eccentricity relative to the selected
demand set:

\[
E(v)
=
\max_iD_{iv}.
\]

The vertex center is

\[
v^*
=
\argmin_{v\in V}E(v).
\]

Computationally, this is a columnwise maximum over \(D\), followed by an
\texttt{argmin}.

The resulting radius

\[
R_V
=
\min_{v\in V}\max_iD_{iv}
\]

also provides a valid incumbent for the continuous absolute-center search.

On trees, the classical graph center goes back to Jordan
\cite{Jordan1869}.

\section{Absolute 1-center}

The absolute-center problem requires searching edge interiors.

Consider edge

\[
e=(u,w)
\]

with length \(L\).

Parameterize a point \(x(t)\) on the edge by distance

\[
t\in[0,L]
\]

from endpoint \(u\).

For demand \(q_i\), define

\[
a_i=d_G(q_i,u),
\qquad
b_i=d_G(q_i,w).
\]

Any route from the interior point to the remainder of the graph must leave the
edge through one of its endpoints.

Therefore

\begin{equation}
d_G(q_i,x(t))
=
\min\{a_i+t,\;b_i+L-t\}.
\label{eq:tent}
\end{equation}

This is a tent-shaped piecewise-linear function.

\subsection{Switch point}

The two endpoint routes are equal when

\[
a_i+t
=
b_i+L-t.
\]

Thus the switch point is

\begin{equation}
t_i^*
=
\frac{b_i-a_i+L}{2}.
\label{eq:switch}
\end{equation}

Because \(u\) and \(w\) are connected by an edge of length \(L\),

\[
|a_i-b_i|\le L
\]

by the triangle inequality, hence

\[
t_i^*\in[0,L]
\]

in exact arithmetic.

\subsection{Upper envelope}

The eccentricity of \(x(t)\) is

\[
E_e(t)
=
\max_i d_G(q_i,x(t)).
\]

Sort the switch points:

\[
t_{(1)}^*
\le
t_{(2)}^*
\le
\cdots
\le
t_{(k)}^*.
\]

Between consecutive switch points, the sets of demands using the two endpoint
branches remain fixed.

On one such interval,

\[
E_e(t)
=
\max\{A+t,\;B+L-t\},
\]

where \(A\) is the largest \(a_i\) among demands still using the \(u\)-side
branch, and \(B\) is the largest \(b_i\) among demands using the \(w\)-side
branch.

The first term rises with slope \(+1\) and the second falls with slope \(-1\).

Their unconstrained intersection occurs at

\[
\widehat t
=
\frac{B+L-A}{2}.
\]

For interval

\[
[\underline t,\overline t],
\]

the interval minimizer is

\[
t^*
=
\clip(
\widehat t;
\underline t,
\overline t
).
\]

Evaluating all breakpoint intervals yields the exact minimum on the edge,
modulo floating-point precision.

\subsection{Prefix and suffix maxima}

After sorting the switch points, the implementation computes:

\begin{itemize}[leftmargin=1.5em]

\item suffix maxima of the \(a_i\);

\item prefix maxima of the \(b_i\).

\end{itemize}

This gives all values of \(A\) and \(B\) in linear work after sorting.

The resulting per-edge complexity is

\[
O(k\log k),
\]

dominated by sorting.

This is a direct computational realization of the classical absolute-center
edge structure associated with Hakimi and Kariv--Hakimi
\cite{Hakimi1964,KarivHakimi1979a}, with textbook treatment in Handler and
Mirchandani \cite{HandlerMirchandani1979} and Daskin
\cite{Daskin2013}.

The particular NumPy prefix/suffix organization is an implementation choice,
not a separate location algorithm.

\subsection{Related literature not implemented}

Minieka \cite{Minieka1970} studied the related \(m\)-center problem and
continuous-network variants.

Megiddo \cite{Megiddo1983} later developed parametric-search techniques that
remove logarithmic factors in related network-center computations. Those
refinements are not implemented here.

\section{Safe edge pruning}

Sweeping every edge is often unnecessary.

The vertex-center radius

\[
R_V
=
\min_{v\in V}\max_iD_{iv}
\]

is already a feasible absolute-center solution.

For any point \(x(t)\) on edge \((u,w)\),

\[
d_G(q_i,x(t))
=
\min\{D_{iu}+t,\;D_{iw}+L-t\}.
\]

Because \(t\ge0\) and \(L-t\ge0\),

\[
d_G(q_i,x(t))
\ge
\min\{D_{iu},D_{iw}\}.
\]

Therefore every point on the edge satisfies

\begin{equation}
E_e(t)
\ge
LB_e
:=
\max_i\min\{D_{iu},D_{iw}\}.
\label{eq:lowerbound}
\end{equation}

If \(LB_e\) cannot improve the incumbent radius, the edge need not be swept.

This is exact branch-and-bound logic in real arithmetic.

It is an implementation acceleration rather than a different facility-location
objective.

\section{Numerical safeguards}

The mathematical algorithms operate over exact real-valued network distances,
while the implementation uses floating-point arithmetic.

Several safeguards are therefore applied.

\subsection{Breakpoint clipping}

The triangle inequality should place every switch point in \([0,L]\).

Tiny numerical violations are defensively clipped:

\[
t_i^*
\leftarrow
\clip(t_i^*;0,L).
\]

\subsection{Pruning tolerance}

When the distance matrix uses reduced precision, a lower bound may differ
slightly from its exact value.

A dtype- and scale-aware margin is included when comparing lower bounds with the
incumbent radius so a candidate edge is not discarded solely because of
floating-point rounding.

\subsection{Meaning of exactness}

The absolute-center sweep is algorithmically exact for the supplied
shortest-path matrix.

Thus, ``exact'' means:

\begin{quote}
exact for the stored network metric, modulo explicit floating-point tolerance.
\end{quote}

A float32 run is therefore exact with respect to the float32 distance matrix,
but that matrix is itself a reduced-precision approximation to the
corresponding float64 metric.

\section{Parallel execution and memory management}

The package separates mathematical computation from execution strategy.

Parallelism does not change the mathematical objective or intended result.

\subsection{Shortest-path blocks}

Demand sources are divided into bounded blocks.

Each block is passed to SciPy's Dijkstra routine.

Results are written directly into a preallocated final matrix.

\subsection{Absolute-center blocks}

The edge sweep operates on bounded groups of segment--demand cells.

The \texttt{max\_cells} parameter controls peak temporary storage.

\subsection{Worker policy}

One worker is the conservative default.

Additional workers may be requested explicitly.

The package exposes process- and thread-based backends because their
performance and memory trade-offs depend on:

\begin{itemize}[leftmargin=1.5em]
\item graph size;
\item demand count;
\item operating system;
\item SciPy build;
\item available RAM;
\item memory bandwidth.
\end{itemize}

No backend is assumed to be universally superior.

\section{Complete computational pipeline}

The first public release can be summarized as follows.

\begin{lstlisting}
BUILD-NETWORK(linework, topology options):

    project linework to a metre-based CRS
    remove unusable geometries

    if planar noding requested:
        split at geometric crossings

    else if shared-vertex noding enabled:
        split where source vertices are shared
        by distinct input lines

    preserve closed rings
    identify network nodes
    remove slivers and self-loops
    compact unused node IDs

    create undirected sparse adjacency
    retain minimum direct weight for parallel node pairs

    optionally restrict to a connected component

    return network


PREPARE-DEMAND(demand):

    convert non-point geometries to representative points
    transform to network CRS
    snap each observation to nearest network node
    enforce optional maximum snap distance
    coalesce duplicate snapped nodes

    return unique demand nodes and median weights


SOLVE(network, demand):

    D <- shortest-path distances from demand nodes
         to all network vertices using Dijkstra

    median <-
        argmin_v sum_i w_i D[i,v]

    vertex_center <-
        argmin_v max_i D[i,v]

    incumbent <- vertex_center

    for each edge e=(u,w,L):

        LB[e] <-
            max_i min(D[i,u], D[i,w])

    discard edges that cannot improve incumbent

    for each surviving edge:

        a_i <- D[i,u]
        b_i <- D[i,w]

        switch_i <-
            (b_i - a_i + L) / 2

        sort demands by switch_i

        compute suffix maxima of a_i
        compute prefix maxima of b_i

        analytically minimize
            max(A + t, B + L - t)
        on every breakpoint interval

    absolute_center <-
        best of vertex incumbent
        and surviving edge-interior candidates

    return median,
           vertex_center,
           absolute_center
\end{lstlisting}

\section{Complexity}

A useful high-level complexity summary is:

\begin{center}
\begin{tabular}{p{0.30\linewidth}p{0.26\linewidth}p{0.35\linewidth}}
\toprule
\textbf{Stage} & \textbf{Indicative cost} & \textbf{Source / comment} \\
\midrule

Planar intersection reporting
&
\(O((s+i)\log s)\)
&
Bentley--Ottmann for \(s\) segments and \(i\) intersections.
\\

Nearest-node demand mapping
&
approximately \(O(k\log n)\)
&
k-d tree; behavior depends on data distribution.
\\

Shortest-path matrix
&
approximately \(O(k(m+n\log n))\)
&
Dijkstra; often the dominant stage.
\\

Weighted 1-median
&
\(O(kn)\)
&
Hakimi vertex reduction followed by weighted column sum.
\\

Vertex 1-center
&
\(O(kn)\)
&
Columnwise maximum followed by minimum.
\\

Edge lower bounds
&
\(O(km)\)
&
Blocked and vectorized.
\\

Absolute-center sweep
&
\(O(mk\log k)\) worst case
&
Only unpruned edges are swept in practice.
\\

\bottomrule
\end{tabular}
\end{center}

The shortest-path matrix requires \(O(kn)\) memory and is frequently the
practical bottleneck.

\section{Implementation work beyond the published algorithms}

The mathematical objectives and their classical structure are not novel.

The following implementation choices were developed during pre-release work,
including AI-assisted development sessions. They should be understood as
engineering on top of the published algorithms rather than as new facility
location methods.

\subsection{Optimizations}

\begin{enumerate}[leftmargin=1.5em]

\item \textbf{Segment pruning by a lower bound.}

The bound

\[
LB_e
=
\max_i \min\{D_{iu},D_{iw}\}
\]

is used to eliminate edges that cannot improve the best known vertex radius.

\item \textbf{Batched vectorization of the sweep.}

Many candidate edges are solved simultaneously using NumPy arrays. Prefix and
suffix running maxima replace repeated interval-by-interval loops.

\item \textbf{Memory-budgeted sweep blocking.}

The absolute-center sweep processes bounded groups of segment--demand cells so
temporary memory does not scale directly with the entire edge set.

\item \textbf{Contiguous sort inputs.}

Endpoint-distance arrays are copied into row-contiguous blocks before
per-edge sorting to simplify and accelerate vectorized operations.

\item \textbf{Conservative parallelism.}

Worker count is treated as a maximum, not a promise. Small workloads remain
serial when worker startup would dominate useful work.

\item \textbf{Bounded shortest-path chunks.}

Dijkstra source nodes are processed in bounded row groups and cast promptly to
the requested storage dtype.

\item \textbf{Streaming parallel assembly.}

Parallel Dijkstra results are consumed in order and written directly into the
final matrix rather than accumulated and stacked afterward.

\item \textbf{Exact demand coalescing.}

Repeated observations snapped to the same node are solved once. Counts or
weights are preserved exactly for the median, while duplicates vanish from the
minimax objective.

\item \textbf{Scale-aware pruning tolerance.}

The pruning margin widens according to the stored floating-point dtype so a
borderline edge is not removed merely because of rounding.

\end{enumerate}

\subsection{An optimization tried and rejected}

A best-first bound-tightening strategy was explored during pre-release
development.

The idea was to sweep the most promising edges first, update the incumbent
radius, and then repeat pruning with the tighter bound.

This is standard branch-and-bound practice.

It was removed because measurements showed that the first lower-bound pass
already reduced the candidate edge set sufficiently that the second pruning
stage did not justify the added complexity.

This history is retained to prevent future optimization work from repeating the
same experiment without measurement.

\section{Correctness bugs found and fixed during pre-release development}

Several important failures were identified before the first public release.
These are worth documenting because they could produce plausible wrong answers
rather than obvious crashes.

\subsection{Parallel edges silently summed}

Sparse COO-to-CSR conversion can add duplicate entries.

If two roads join the same pair of graph nodes and their lengths are summed,
the resulting direct edge weight may become larger than a valid alternative
path.

This can violate the metric structure expected by the absolute-center sweep.

The current implementation collapses parallel edges using the minimum direct
length.

\subsection{Closed loops deleted}

A closed road may merge into a LineString whose endpoints coincide.

If the result is then interpreted as a self-loop and removed, the entire ring
disappears from the graph.

The current graph builder splits such rings before self-loop filtering.

\subsection{Empty geometries corrupting endpoint lookup}

Empty geometries can produce invalid coordinate runs and incorrect endpoint
association if not removed before indexing.

The graph builder filters unusable geometries before endpoint extraction.

\subsection{Needless reprojection}

Automatically reprojecting all input to a locally inferred UTM zone is
appropriate for geographic coordinates but can be wrong for input already
projected in valid metre units.

The current implementation preserves projected metre-based input and only
reprojects when necessary.

\subsection{Nearly closed rings deleted}

A ring can be visually closed while its first and last coordinates differ by
floating-point noise.

Exact closure tests may miss this case.

The graph builder uses the same snapping tolerance used for node identity when
detecting nearly closed rings.

\subsection{False junctions at grade-separated crossings}

Treating every two-dimensional line crossing as connected can create impossible
turns between bridges, flyovers, tunnels, and roads below them.

Full planar noding is therefore explicit rather than the default.

\subsection{Shared-vertex noding initially too narrow}

An earlier pre-release rule handled T-junctions where one line ended on the
interior of another but still missed cases where two lines both continued
through a shared source vertex.

The current rule defines a junction by participation of at least two distinct
input LineStrings at the snapped source coordinate.

\subsection{Filtered slivers leaving orphan graph nodes}

Assigning node IDs before removing slivers and snapped self-loops can leave
unused rows and columns in the sparse graph.

The current implementation compacts node IDs after edge filtering.

\section{Validation strategy}

Validation separates facility-location mathematics from GIS graph-construction
correctness as much as possible.

\subsection{Hand-computable cases}

Small networks with known answers test:

\begin{itemize}[leftmargin=1.5em]
\item weighted median locations;
\item vertex centers;
\item edge-interior absolute centers.
\end{itemize}

\subsection{Independent edge sampling}

The continuous absolute-center solution is compared with independent brute-force
sampling along network edges on randomly generated test networks.

The sampled solution is approximate, but it shares no breakpoint logic with the
exact sweep.

It therefore provides an independent check of the sweep implementation.

\subsection{What brute force does not validate}

Brute-force edge sampling validates the optimizer conditional on the supplied
graph and shortest-path matrix.

It does not validate whether the graph correctly represents the intended road
system.

For example, no numerical test can determine from geometry alone whether a
bridge should connect to a road below it.

That remains a topology and data-semantics question.

\subsection{Regression coverage}

Regression tests cover failures including:

\begin{itemize}[leftmargin=1.5em]

\item parallel edges being combined incorrectly;

\item closed rings disappearing as self-loops;

\item nearly closed rings being lost because of floating-point differences;

\item empty geometries corrupting node assignment;

\item false connections at geometric bridge crossings;

\item missed T-junctions;

\item missed shared interior-interior junctions;

\item filtered edges leaving isolated graph nodes;

\item invalid projected units;

\item disconnected demand/network combinations;

\item worker-count and block-size invariance;

\item float32 pruning behavior;

\item package-version and CLI consistency.

\end{itemize}

\section{Pre-release development history}

Before v0.1.0, the codebase underwent several internal development and review
iterations.

Those intermediate states were not intended to define a public semantic-version
sequence and are therefore described here simply as \emph{pre-release
development history}.

The important distinction is between:

\begin{enumerate}[leftmargin=1.5em]

\item changes to the mathematical facility-location problem; and

\item changes to the graph, numerical representation, and implementation on
which that problem is solved.

\end{enumerate}

The core optimization problems remained the same throughout pre-release
development:

\begin{itemize}[leftmargin=1.5em]

\item weighted vertex 1-median;

\item vertex-restricted minimax center;

\item continuous absolute 1-center.

\end{itemize}

Most substantive pre-release changes instead concerned:

\begin{itemize}[leftmargin=1.5em]

\item safer road-network topology;

\item correct treatment of parallel edges;

\item preservation of rings;

\item coordinate-system handling;

\item duplicate-demand reduction;

\item bounded-memory shortest paths;

\item numerical validation;

\item parallel execution;

\item regression coverage.

\end{itemize}

These changes can alter a numerical answer when they alter the graph \(G\) being
represented, but they do not define a different median or center objective.

\section{Audit: same mathematics versus changed preprocessing}

The table below summarizes the distinction.

\small

\begin{longtable}{p{0.21\linewidth}p{0.29\linewidth}p{0.29\linewidth}p{0.15\linewidth}}
\toprule
\textbf{Aspect}
&
\textbf{Early prototype}
&
\textbf{v0.1.0}
&
\textbf{Changed mathematics?}
\\
\midrule
\endfirsthead

\toprule
\textbf{Aspect}
&
\textbf{Early prototype}
&
\textbf{v0.1.0}
&
\textbf{Changed mathematics?}
\\
\midrule
\endhead

Weighted 1-median
&
Weighted total shortest-path distance over vertices.
&
Same objective and vertex reduction.
&
No.
\\
\addlinespace

Vertex 1-center
&
Maximum demand distance minimized over vertices.
&
Same objective and computation.
&
No.
\\
\addlinespace

Absolute 1-center
&
Continuous edge-interior minimax problem using breakpoint structure.
&
Same mathematical problem and same edgewise envelope structure.
&
No.
\\
\addlinespace

Shortest paths
&
Dijkstra on an undirected sparse graph.
&
Same shortest-path metric, with stronger validation and bounded-memory execution.
&
No.
\\
\addlinespace

Edge pruning
&
Lower bound based on endpoint distances.
&
Same bound with more conservative floating-point handling.
&
No.
\\
\addlinespace

Demand duplicates
&
Could produce repeated shortest-path rows.
&
Coalesced exactly before shortest-path computation.
&
No.
\\
\addlinespace

Crossing treatment
&
Earlier prototype could planarize all crossings.
&
Shared-source-vertex noding is default; full planar noding is explicit.
&
Preprocessing changed.
\\
\addlinespace

Shared interior junctions
&
Not fully recovered in the earliest rule.
&
Recovered when distinct source LineStrings share a snapped source vertex.
&
Preprocessing changed.
\\
\addlinespace

Connected component choice
&
Could depend on node count.
&
Uses total road length.
&
Preprocessing changed.
\\
\addlinespace

CRS handling
&
Weaker assumptions about projected units.
&
Metre-based projected CRS required or constructed explicitly.
&
Measurement policy changed.
\\
\addlinespace

Parallel-edge handling
&
Duplicate sparse entries could sum.
&
Minimum direct edge length retained.
&
Bug fix.
\\
\addlinespace

Closed rings
&
Could disappear as self-loops.
&
Preserved by splitting before graph construction.
&
Bug fix.
\\
\addlinespace

Unused node IDs
&
Could survive edge filtering.
&
Compacted after filtering.
&
Bug fix.
\\
\addlinespace

Precision
&
Primarily float64-oriented.
&
float32 storage optional with safe pruning tolerance.
&
Optional numerical approximation only.
\\

\bottomrule
\end{longtable}

\normalsize

The most important point is that the location objective and the classical
absolute-center method remained stable. The pre-release changes mostly improved
the graph to which those algorithms are applied.

\section{Scope and limitations}

The first public release deliberately addresses a restricted class of
network-location problems.

\subsection{Undirected networks}

The graph is undirected.

The package does not currently model:

\begin{itemize}[leftmargin=1.5em]
\item one-way streets;
\item asymmetric travel costs;
\item turn restrictions;
\item time-dependent routing.
\end{itemize}

These require a different state representation and, in particular, a
reconsideration of continuous edge-interior treatment.

\subsection{Vertex demand}

External demand is snapped to graph vertices.

Demand is not currently represented continuously along edge interiors.

\subsection{Unweighted minimax objective}

Weights affect the 1-median but not the center.

The current center objective is

\[
\min_x\max_i d_G(x,q_i),
\]

not a weighted minimax variant.

A weighted minimax objective changes the edgewise functions and cannot simply
reuse the current prefix/suffix structure unchanged.

\subsection{One facility}

The package solves the \(p=1\) case.

General \(p\)-center and \(p\)-median problems are substantially more difficult
and belong to a broader optimization literature
\cite{KarivHakimi1979a,KarivHakimi1979b}.

\subsection{Static edge costs}

Edge costs represent a fixed metric.

Congestion, traffic variation, schedules, stochastic travel time, and other
dynamic costs are outside the present scope.

\subsection{Topology semantics}

Geometry alone cannot fully determine real transportation connectivity.

The graph builder does not automatically infer every bridge, tunnel, layer,
access, one-way, or turn-restriction semantic from external road-database
attributes.

The analyst remains responsible for ensuring that the input network represents
the intended movement system.

\section{Network median versus geographic centroid}

The network median should not be confused with a Euclidean centroid.

For Euclidean coordinates \(x_i\), a weighted arithmetic centroid is

\[
\bar{x}
=
\frac{\sum_i w_i x_i}
{\sum_i w_i}.
\]

Equivalently, it minimizes a squared Euclidean-distance objective.

The network median instead minimizes

\[
\sum_i w_i d_G(x,q_i).
\]

It is therefore more naturally interpreted as a \emph{minisum accessibility
center} than as a literal coordinate centroid.

A closer network analogue of a Euclidean squared-distance centroid would be a
Fr\'echet mean or network barycenter:

\[
x^*
=
\argmin_{x\in G}
\sum_i w_i d_G(x,q_i)^2.
\]

That problem is not implemented in v0.1.0.

\section{Implementation versus mathematical contribution}

It is important to distinguish published facility-location mathematics from the
package implementation.

The following are classical:

\begin{itemize}[leftmargin=1.5em]

\item the network median objective;

\item the absolute-center objective;

\item vertex optimality of the network median for vertex demand;

\item edge-interior absolute centers;

\item shortest-path distance as the network metric;

\item the breakpoint structure of the continuous center problem.

\end{itemize}

The following are implementation and engineering choices in \texttt{netcenter}:

\begin{itemize}[leftmargin=1.5em]

\item GIS linework preprocessing;

\item shared-source-vertex topology recovery;

\item defensive handling of rings, slivers, and parallel edges;

\item demand coalescing;

\item NumPy vectorization;

\item prefix/suffix envelope evaluation;

\item lower-bound pruning;

\item blocked shortest-path computation;

\item streaming parallel assembly;

\item reduced-precision storage support;

\item numerical tolerance policy;

\item regression-test coverage.

\end{itemize}

No novelty claim is made for the underlying network-location algorithms.

\section{Conclusion}

\texttt{netcenter} v0.1.0 is best understood as a practical geospatial
implementation of classical network-location theory.

Its mathematical core is compact:

\begin{enumerate}[leftmargin=1.5em]

\item construct a metric graph from linework;

\item map demand to network vertices;

\item compute shortest-path distances;

\item solve the weighted 1-median directly over vertices;

\item solve the vertex minimax center directly over vertices;

\item use that center as an incumbent;

\item prune edges whose lower bounds cannot improve the incumbent;

\item solve the remaining edge-interior minimax problems analytically through
breakpoint sorting and upper-envelope minimization.

\end{enumerate}

The principal difficulty in practical use is not defining the objective but
ensuring that the graph faithfully represents the intended transportation
network.

Accordingly, much of the package's engineering is devoted to topology,
coordinate systems, memory safety, numerical safeguards, provenance, and
validation rather than to modifying the classical facility-location
mathematics.

\begin{thebibliography}{99}

\bibitem{Bavelas1950}
A. Bavelas.
\newblock Communication patterns in task-oriented groups.
\newblock \emph{Journal of the Acoustical Society of America},
22(6):725--730, 1950.

\bibitem{Bentley1975}
J. L. Bentley.
\newblock Multidimensional binary search trees used for associative searching.
\newblock \emph{Communications of the ACM}, 18(9):509--517, 1975.

\bibitem{BentleyOttmann1979}
J. L. Bentley and T. A. Ottmann.
\newblock Algorithms for reporting and counting geometric intersections.
\newblock \emph{IEEE Transactions on Computers}, C-28(9):643--647, 1979.

\bibitem{Daskin2013}
M. S. Daskin.
\newblock \emph{Network and Discrete Location: Models, Algorithms, and Applications}.
\newblock 2nd edition, Wiley, 2013.
\newblock doi:10.1002/9781118537015.

\bibitem{Dijkstra1959}
E. W. Dijkstra.
\newblock A note on two problems in connexion with graphs.
\newblock \emph{Numerische Mathematik}, 1:269--271, 1959.

\bibitem{Freeman1978}
L. C. Freeman.
\newblock Centrality in social networks: conceptual clarification.
\newblock \emph{Social Networks}, 1(3):215--239, 1978.

\bibitem{Goldman1971}
A. J. Goldman.
\newblock Optimal center location in simple networks.
\newblock \emph{Transportation Science}, 5(2):212--221, 1971.

\bibitem{GreeneYao1986}
D. H. Greene and F. F. Yao.
\newblock Finite-resolution computational geometry.
\newblock In \emph{27th Annual Symposium on Foundations of Computer Science},
pages 143--152, 1986.

\bibitem{Hakimi1964}
S. L. Hakimi.
\newblock Optimum locations of switching centers and the absolute centers and
medians of a graph.
\newblock \emph{Operations Research}, 12(3):450--459, 1964.
\newblock doi:10.1287/opre.12.3.450.

\bibitem{Hakimi1965}
S. L. Hakimi.
\newblock Optimum distribution of switching centers in a communication network
and some related graph theoretic problems.
\newblock \emph{Operations Research}, 13(3):462--475, 1965.
\newblock doi:10.1287/opre.13.3.462.

\bibitem{HandlerMirchandani1979}
G. Y. Handler and P. B. Mirchandani.
\newblock \emph{Location on Networks: Theory and Algorithms}.
\newblock MIT Press, Cambridge, Massachusetts, 1979.

\bibitem{Hobby1999}
J. D. Hobby.
\newblock Practical segment intersection with finite precision output.
\newblock \emph{Computational Geometry}, 13(4):199--214, 1999.

\bibitem{HopcroftTarjan1973}
J. Hopcroft and R. Tarjan.
\newblock Algorithm 447: efficient algorithms for graph manipulation.
\newblock \emph{Communications of the ACM}, 16(6):372--378, 1973.

\bibitem{Jordan1869}
C. Jordan.
\newblock Sur les assemblages de lignes.
\newblock \emph{Journal f\"ur die reine und angewandte Mathematik},
70:185--190, 1869.

\bibitem{KarivHakimi1979a}
O. Kariv and S. L. Hakimi.
\newblock An algorithmic approach to network location problems.
I: The \(p\)-centers.
\newblock \emph{SIAM Journal on Applied Mathematics},
37(3):513--538, 1979.
\newblock doi:10.1137/0137040.

\bibitem{KarivHakimi1979b}
O. Kariv and S. L. Hakimi.
\newblock An algorithmic approach to network location problems.
II: The \(p\)-medians.
\newblock \emph{SIAM Journal on Applied Mathematics},
37(3):539--560, 1979.
\newblock doi:10.1137/0137041.

\bibitem{Megiddo1983}
N. Megiddo.
\newblock Applying parallel computation algorithms in the design of serial
algorithms.
\newblock \emph{Journal of the ACM}, 30(4):852--865, 1983.

\bibitem{Minieka1970}
E. Minieka.
\newblock The \(m\)-center problem.
\newblock \emph{SIAM Review}, 12(1):138--139, 1970.

\bibitem{MirchandaniFrancis1990}
P. B. Mirchandani and R. L. Francis, editors.
\newblock \emph{Discrete Location Theory}.
\newblock Wiley, 1990.

\bibitem{SciPyDijkstra}
SciPy developers.
\newblock \texttt{scipy.sparse.csgraph.dijkstra} documentation.
\newblock \url{https://docs.scipy.org/doc/scipy/reference/generated/scipy.sparse.csgraph.dijkstra.html}.

\bibitem{ShapelyNode}
Shapely developers.
\newblock \texttt{shapely.node} documentation.
\newblock \url{https://shapely.readthedocs.io/en/stable/reference/shapely.node.html}.

\end{thebibliography}

\end{document}
